% Демо-презентація етапу 4 (структура — docs/stage4.md). Цифри — з pnpm test:report і OpenAPI.
\documentclass[aspectratio=169,11pt]{beamer}
\usetheme{metropolis}
\usepackage{fontspec}
\usepackage{polyglossia}
\setmainlanguage{ukrainian}
\setmainfont{DejaVuSans}[Extension=.ttf,UprightFont=*,BoldFont=*-Bold,ItalicFont=*-Oblique,BoldItalicFont=*-BoldOblique]
\setsansfont{DejaVuSans}[Extension=.ttf,UprightFont=*,BoldFont=*-Bold,ItalicFont=*-Oblique,BoldItalicFont=*-BoldOblique]
\setmonofont{DejaVuSansMono}[Extension=.ttf,UprightFont=*,BoldFont=*-Bold,ItalicFont=*-Oblique,BoldItalicFont=*-BoldOblique,Scale=0.9]
\usepackage{booktabs,tabularx,array,colortbl,xcolor}
\definecolor{sgreen}{HTML}{E2F0D9}
\definecolor{wred}{HTML}{FBE5D6}
\definecolor{oblue}{HTML}{DEEBF7}
\definecolor{tyellow}{HTML}{FFF2CC}
\definecolor{ink}{HTML}{2132A8}
\newcolumntype{Y}{>{\raggedright\arraybackslash}X}
\setbeamerfont{frametitle}{size=\large}
\setbeamertemplate{itemize items}{\textbullet}
\setlength{\leftmargini}{1.2em}
\graphicspath{{img/}{../docs/uml/out/}}
\newcommand{\todo}[1]{\textcolor{red}{\textbf{[TODO: #1]}}}
\newcommand{\img}[2][0.95\textwidth]{\IfFileExists{img/#2}{\includegraphics[width=#1,height=0.78\textheight,keepaspectratio]{#2}}{\IfFileExists{../docs/uml/out/#2}{\includegraphics[width=#1,height=0.78\textheight,keepaspectratio]{#2}}{\todo{#2}}}}
\newcommand{\bignum}[2]{\begin{minipage}[t]{0.23\textwidth}\centering{\Huge\bfseries\color{ink}#1}\\[2pt]{\small #2}\end{minipage}}

\title{StudyFlow: реалізація прототипу}
\subtitle{Інформаційні технології, Модуль 2 — Етап 4}
\author{Петров Богдан \and Мекшаков Артемій}
\institute{Група МІ-41}
\date{18 листопада 2026}

\begin{document}
\maketitle

\begin{frame}{StudyFlow в одному слайді}
\footnotesize
\begin{columns}[T]
\column{0.5\textwidth}
\textbf{Проблема}
\begin{itemize}
  \item дедлайни розкидані між Classroom і чатами
  \item староста дублює їх вручну
  \item не видно, хто вже здав
\end{itemize}
\textbf{Рішення}
\begin{itemize}
  \item староста створює дедлайн один раз --- він у календарі кожного
  \item кожен відмічає виконання сам
  \item група бачить прогрес у відсотках
\end{itemize}
\column{0.46\textwidth}
\textbf{Стек}
\begin{itemize}
  \item Next.js 16, Tailwind 4, shadcn/ui, TanStack Query
  \item NestJS 12, Prisma 7, PostgreSQL 16
  \item JWT (access у пам'яті, refresh у httpOnly cookie)
  \item Vitest, supertest, Playwright
  \item Docker Compose, GitHub Actions, Swagger
\end{itemize}
\end{columns}
\end{frame}

\begin{frame}{План спринта: план vs факт}
\begin{columns}[T]
\column{0.44\textwidth}
\small
\begin{tabular}{@{}lccc@{}}
\toprule
Спринт & ID & План SP & Done SP\\
\midrule
1 & 11 & 22 & 22\\
2 & 11 & 24 & 24\\
3 & 9 & 22 & 22\\
\midrule
\textbf{Разом} & \textbf{31} & \textbf{68} & \textbf{68}\\
\bottomrule
\end{tabular}
\medskip

Усі 31 елемент спринтів 1--3 --- Done: Must, Should і Could (US-21, перетягування).
\column{0.54\textwidth}
\centering\img{notion-sprints.png}
\end{columns}
\end{frame}

\begin{frame}{Архітектура: модульний моноліт}
\centering\img[0.92\textwidth]{component.png}
\end{frame}

\begin{frame}{Модель даних: 8 таблиць}
\centering\img[0.8\textwidth]{er.png}
\end{frame}

\begin{frame}{UML: варіанти використання і груповий дедлайн}
\centering\img[0.3\textwidth]{use-case.png}\hfill\img[0.66\textwidth]{sequence-group-deadline.png}
\end{frame}

\begin{frame}{Демо поточної версії}
\begin{enumerate}
  \item Староста: група «МІ-41» $\to$ новий груповий дедлайн на завтра з предметом
  \item Учасник (інкогніто): дедлайн уже в календарі тижня $\to$ відмітити виконаним
  \item Староста: оновити сторінку групи --- прогрес зріс
  \item Канбан: перетягнути задачу «В роботі» $\to$ «Готово»
  \item Swagger: \texttt{POST /api/dev/reminders/run} $\to$ лист у Mailpit
  \item Swagger \texttt{/api/docs}: теги і схеми DTO
\end{enumerate}
\end{frame}

\begin{frame}{Скріншоти: дедлайни і календар}
\centering\img[0.49\textwidth]{dashboard.png}\hfill\img[0.49\textwidth]{calendar-week.png}
\end{frame}

\begin{frame}{Скріншоти: група, канбан, нагадування}
\centering\img[0.325\textwidth]{group.png}\hfill\img[0.325\textwidth]{board.png}\hfill\img[0.325\textwidth]{mailpit.png}
\end{frame}

\begin{frame}{Модульні тести}
\centering
\bignum{40/40}{пройдено}\hfill\bignum{54\,\%}{рядків покрито}\hfill\bignum{8}{модулів з тестами}\hfill\bignum{1{,}2 с}{прогін}
\vspace{1.2em}
\raggedright\small
\begin{itemize}
  \item Vitest, Prisma замокано (\texttt{vitest-mock-extended}), сервіси без DI
  \item UT-07 refresh зі старим токеном $\to$ 401 \quad UT-15 останнього старосту не понизити
  \item UT-27 2 з 4 $\to$ 50\,\% \quad UT-33 повторний запуск нагадувань не шле дублікатів
  \item журнал з ID --- у Додатку~А звіту (\texttt{pnpm test:report})
\end{itemize}
\end{frame}

\begin{frame}{Інтеграційні та E2E-тести}
\footnotesize
\begin{tabularx}{\textwidth}{@{}l Y c@{}}
\toprule
Набір & Що перевіряє & Результат\\
\midrule
Інтеграційні, 13 & реальна БД \texttt{studyflow\_test}, HTTP через supertest, пошта --- заглушка: auth і cookie, ролі (403), видимість дедлайнів, прогрес 50\,\%, фільтри, задачі, нагадування без дублікатів, Swagger & 13/13\\
\midrule
E2E, 3 & Playwright, два браузери: реєстрація $\to$ група $\to$ груповий дедлайн $\to$ вступ за кодом $\to$ відмітка $\to$ 50\,\%; невірний пароль; дедлайн у правильний день календаря & 3/3\\
\bottomrule
\end{tabularx}
\medskip

Покриття коду інтеграційними тестами --- 73{,}5\,\% рядків. Інтеграційні тести входять у CI (GitHub Actions + PostgreSQL), прогін 11.11 --- зелений.
\end{frame}

\begin{frame}{API-документація}
\begin{columns}[T]
\column{0.36\textwidth}
\small
\begin{itemize}
  \item 34 операції, 8 тегів
  \item Swagger UI --- \texttt{/api/docs}
  \item \texttt{docs/api/openapi.json} у репозиторії
  \item з нього генеруються типи фронтенду
  \item помилки й DTO з прикладами, Bearer або cookie
\end{itemize}
\column{0.62\textwidth}
\centering\img{swagger.png}
\end{columns}
\end{frame}

\begin{frame}{Рефлексія спринта}
\footnotesize
\renewcommand{\arraystretch}{1.0}
\begin{tabularx}{\textwidth}{|Y|Y|}
\hline
\cellcolor{sgreen}\textbf{Успіхи} & \cellcolor{wred}\textbf{Неуспіхи}\\
\hline
\cellcolor{sgreen!40}\begin{itemize}
  \item 31 з 31 елемента, 68 SP --- включно з Should і Could
  \item тестів більше за план: 40 / 13 / 3
  \item правило $\to$ тест $\to$ журнал за ID
\end{itemize} &
\cellcolor{wred!40}\begin{itemize}
  \item реалізація стартувала 10.11 замість 28.10
  \item коміти з одного акаунта, без PR і рев'ю
  \item CI у GitHub --- лише з 11.11
\end{itemize}\\
\hline
\cellcolor{oblue}\textbf{Причини} & \cellcolor{tyellow}\textbf{Дії на спринт 4}\\
\hline
\cellcolor{oblue!40}\begin{itemize}
  \item спринти 1--2 пішли на документацію етапів 2--3
  \item дедлайн етапу 4 --- 15.11 замість 18.11
\end{itemize} &
\cellcolor{tyellow!40}\begin{itemize}
  \item feature-гілки з власних акаунтів, PR з рев'ю, CI
  \item деплой першим; статуси в Notion після злиття
  \item покриття модульними $\geq$70\,\%
\end{itemize}\\
\hline
\end{tabularx}
\end{frame}

\begin{frame}{Наступні кроки до 01.12}
\begin{itemize}
  \item \textbf{T-09} деплой: web --- Vercel, API --- Railway/Render, керований PostgreSQL; публічне посилання
  \item \textbf{US-24} налаштування нагадувань (поля в БД уже є) і \textbf{US-28} особиста статистика
  \item пілот у групі МІ-41 --- зібрати відгуки
  \item \textbf{D-03} фінальний звіт і презентація модуля
\end{itemize}
\end{frame}

\end{document}
